\chapter{Midterm Review}
\label{ch:review}
\fancyhead[R]{\small\textit{Chapter 7: Midterm Review}}
\section{How to Use This Chapter}
\label{sec:ch7-overview}

Chapters~1--6 built the entire foundational apparatus of finite frame
theory: the definition of a frame, the frame operator and its
eigenvalues, tight frames, dual frames, and the geometric picture
underlying duality. This chapter does not introduce new theory.
Instead, it reorganizes everything you have learned around five central
questions, gives you a single reference page for every major definition
and theorem, and then provides a substantial bank of practice problems
--- a mix of proof-style and computational questions, in the spirit of
the midterm itself.

The material is reorganized \emph{thematically} rather than
chapter-by-chapter, because the five central questions below cut across
the original chapter boundaries, and reviewing this way is a better test
of whether you have internalized the ideas rather than just the order in
which they were presented.

\begin{keyidea}
\textbf{The five central questions of Chapters 1--6:}
\begin{enumerate}[label=(\arabic*)]
  \item \emph{When is a set of vectors a frame?} (Spanning is
        sufficient and necessary in finite dimensions.)
  \item \emph{How do we compute its bounds exactly?} (Eigenvalues of
        the frame operator $S=FF^*$.)
  \item \emph{When is a frame tight, and why does it matter?}
        ($S=AI$; resolution of the identity; roots-of-unity symmetry as
        one mechanism among several.)
  \item \emph{How do we reconstruct a signal from frame coefficients,
        and is the reconstruction unique?} (Always possible via the
        canonical dual; unique only when the frame is a basis;
        otherwise an entire affine family of valid duals exists, and the
        canonical one is optimal in a precise sense.)
  \item \emph{What does that family of duals actually look like?}
        (For the smallest redundant case, an ordinary plane with an
        exact paraboloid of energy over it; the canonical dual sits at
        the bottom, and the shape of the resulting reconstruction-error
        ellipse is governed separately by the frame's own tightness.)
\end{enumerate}
\end{keyidea}

\section{Question 1: When Is a Set of Vectors a Frame?}
\label{sec:review-q1}

\begin{definition}[Finite frame --- Chapter 2]
\label{rev:def-frame}
$\{f_i\}_{i=1}^m\subset\F^n$ is a frame for $\F^n$ if there exist
$0<A\le B<\infty$ with
\[
  A\norm{x}^2 \le \sum_{i=1}^m|\inner{x}{f_i}|^2 \le B\norm{x}^2
  \qquad\forall x\in\F^n.
\]
\end{definition}

\begin{theorem}[Spanning $\Leftrightarrow$ frame --- Chapter 2, Theorem 2.2.1]
\label{rev:thm-spanning}
A finite set $\{f_i\}_{i=1}^m\subset\F^n$ is a frame for $\F^n$ if and
only if it spans $\F^n$.
\end{theorem}

\begin{reviewbox}
This is the single fact that makes finite frame theory tractable:
existence of frame bounds is \emph{automatic} once you check spanning
(equivalently, $\rank F = n$). The upper bound always exists (Cauchy--
Schwarz, $B=\sum_i\norm{f_i}^2$ works); the lower bound is where
spanning is essential (compactness argument: a continuous function on
the unit sphere attains its minimum, and that minimum is forced
positive exactly because no nonzero vector is orthogonal to a spanning
set). \textbf{If an exam question asks ``is this a frame,'' check the
rank of the synthesis matrix --- nothing else is needed.}
\end{reviewbox}

\section{Question 2: How Do We Compute the Bounds Exactly?}
\label{sec:review-q2}

\begin{definition}[Synthesis, analysis, frame operators --- Chapter 3]
\label{rev:def-operators}
$F:\F^m\to\F^n$, $Fc=\sum_ic_if_i$ (matrix with columns $f_i$);
$F^*:\F^n\to\F^m$, $F^*x=(\inner{x}{f_i})_i$ (conjugate transpose);
$S:=FF^*=\sum_if_if_i^*$, the \emph{frame operator}.
\end{definition}

\begin{theorem}[Frame bounds are eigenvalues --- Chapter 3, Theorem 3.5.1]
\label{rev:thm-eigenvalues}
$A_{\mathrm{opt}}=\lambda_{\min}(S)$, $B_{\mathrm{opt}}=\lambda_{\max}(S)$.
\end{theorem}

\begin{reviewbox}
Computationally: \texttt{S = F @ F.T}, then
\texttt{eigh(S)} gives both bounds in one call. \textbf{Always use
\texttt{eigh}, never \texttt{eig}}, since $S$ is symmetric by
construction --- \texttt{eigh} guarantees real, sorted eigenvalues and
an orthonormal eigenbasis, while \texttt{eig} does not.
The proof of Theorem~\ref{rev:thm-eigenvalues} hinges on the
\emph{Rayleigh quotient}: $\inner{Sx}{x}/\norm{x}^2$ is a convex
combination of the eigenvalues of $S$, hence trapped between
$\lambda_{\min}$ and $\lambda_{\max}$, with equality exactly at the
corresponding eigenvectors. This is a proof technique you should be
able to reproduce from scratch, not just cite.
\end{reviewbox}

\begin{proposition}[Trace formula --- Chapter 3, Proposition 3.4.2(iv)]
\label{rev:prop-trace}
$\tr(S) = \sum_{i=1}^m\norm{f_i}^2$. For a tight, unit-norm frame with
$m$ vectors in $\R^n$: $A = m/n$.
\end{proposition}

\section{Question 3: When Is a Frame Tight, and Why?}
\label{sec:review-q3}

\begin{definition}[Tight / Parseval frame --- Chapter 2]
\label{rev:def-tight}
Tight: $A=B$. Equivalently (Chapter 4): $S=AI$. Parseval: tight with
$A=1$.
\end{definition}

\begin{theorem}[Regular polygons --- Chapter 4, Theorem 4.3.1]
\label{rev:thm-polygon}
The vertices of a regular $m$-gon ($m\ge3$) inscribed in the unit
circle form a tight frame for $\R^2$ with $A=m/2$.
\end{theorem}

\begin{theorem}[Projection construction --- Chapter 4, Theorem 4.4.1]
\label{rev:thm-projection}
Projecting any orthonormal basis of $\F^N$ onto a subspace $V$ always
produces a Parseval frame for $V$.
\end{theorem}

\begin{reviewbox}
Two genuinely different \emph{mechanisms} produce tight frames: (a)
\textbf{symmetry} (roots of unity / cyclic group orbits force
trigonometric cross-terms to cancel) and (b) \textbf{projection} (any
orthonormal basis of a bigger space, projected down, is automatically
Parseval, regardless of symmetry). Neither mechanism is the
\emph{definition} of tightness --- tightness is fundamentally an
energy-balance condition, $S=AI$, and these are just two reliable ways
to engineer it. \textbf{Do not write ``tight frames are symmetric'' as
if it were an if-and-only-if statement} --- Proposition 4.5.1 gives an
explicit asymmetric, unequal-norm tight frame.
\end{reviewbox}

\section{Question 4: Reconstruction and Duality}
\label{sec:review-q4}

\begin{definition}[Dual frame --- Chapter 5]
\label{rev:def-dual}
$\{g_i\}$ is dual to $\{f_i\}$ if $x=\sum_i\inner{x}{f_i}g_i$ for all
$x$; equivalently $GF^*=I_n$. The relationship is mutual.
\end{definition}

\begin{proposition}[Canonical dual --- Chapter 5, Proposition 5.3.1]
\label{rev:prop-canonical}
$\tilde f_i := S^{-1}f_i$ is always a valid dual, with $\tilde S=S^{-1}$
and bounds $\tilde A=1/B$, $\tilde B=1/A$.
\end{proposition}

\begin{theorem}[All duals --- Chapter 5, Theorem 5.4.1]
\label{rev:thm-dual-family}
Every valid dual synthesis matrix is $G=\tilde F+Z$ with $ZF^*=0$
(rows of $Z$ in $\ker F$). Unique iff $m=n$ (basis case).
\end{theorem}

\begin{theorem}[Minimal norm --- Chapter 5, Theorem 5.5.1]
\label{rev:thm-minimal-norm}
$\tilde F$ uniquely minimizes $\norm{G}_F^2=\sum_i\norm{g_i}^2$ among
all valid duals, via the Pythagorean decomposition
$\norm{G}_F^2=\norm{\tilde F}_F^2+\norm{Z}_F^2$. Consequently $\tilde F$
also minimizes expected reconstruction error under uncorrelated
coefficient noise.
\end{theorem}

\begin{reviewbox}
\textbf{A commonly confused dimension-counting point, worth getting
exactly right:} for a frame with $m$ vectors in $\F^n$ ($m>n$), the
space of valid perturbations $Z$ (with $ZF^*=0$) has real dimension
$n\cdot(m-n)$ --- \emph{not} $m-n$. The kernel $\ker F\subset\F^m$
itself has dimension $m-n$, but each of the $n$ \emph{rows} of $Z$
ranges independently over that kernel, so the total freedom is $n$
copies of it. For the triangular frame ($m=3$, $n=2$), $\dim\ker F=1$
but the dual family is $2\cdot1=2$-dimensional --- this is exactly the
fact that made the parameter plane of Chapter~6 a genuine \emph{plane}
rather than a line. \textbf{Do not say ``the dual family has dimension
$m-n$''} without the factor of $n$; this is a frequent and easy
mistake.
\end{reviewbox}

\section{Question 5: What Does the Dual Family Look Like?}
\label{sec:review-q5}

\begin{proposition}[The dual family as a plane --- Chapter 6, Proposition 6.2.1]
\label{rev:prop-dual-plane}
For a frame with $m=n+1$ vectors, the affine family of duals is
parametrized by $G(w)=\tilde F+wk^T$, $w\in\R^n$, where $k$ is a unit
vector spanning the (one-dimensional) kernel of $F$.
\end{proposition}

\begin{theorem}[Energy paraboloid --- Chapter 6, Theorem 6.3.1]
\label{rev:thm-paraboloid}
$\norm{G(w)}_F^2=\norm{\tilde F}_F^2+\norm{w}^2$: an exact paraboloid
over the parameter plane, with unique minimum at $w=0$ (the canonical
dual). This holds whether or not the original frame is tight.
\end{theorem}

\begin{proposition}[Error covariance --- Chapter 6, Proposition 6.4.2]
\label{rev:prop-error-cov}
The canonical dual's reconstruction-error covariance under coefficient
noise of variance $\sigma^2$ is $\Sigma_{\tilde F}=\sigma^2S^{-1}$
exactly. Its level sets are circular if and only if the original frame
is tight.
\end{proposition}

\begin{reviewbox}
\textbf{Two different optimization criteria, easy to conflate:} the
canonical dual minimizes \emph{total} energy $\norm{G}_F^2$ over the
\emph{entire} affine family --- always, for any frame, tight or not.
But minimizing total energy is not the same as making the
reconstruction error \emph{isotropic} (equal in every direction): the
shape of the error ellipse at the canonical dual is inherited from
$S^{-1}$, and is only circular if the frame itself was tight to begin
with. Chapter~6's Lab Exercise 6.2 showed explicitly that for a
non-tight frame, there can exist a \emph{different} point on the same
paraboloid where the error ellipse becomes circular --- at the cost of
\emph{more} total energy than the canonical choice. \textbf{Minimal
total error and balanced directional error are, in general, different
points on the paraboloid; the canonical dual only optimizes the
former.}
\end{reviewbox}

\section{Worked Review Problems}
\label{sec:worked-review}

These problems are worked in full, modeling the level of detail expected
on the midterm itself.

\begin{examplebox}[title={Review Problem 1: Is it a frame? Find the bounds.}]
\label{rev:problem1}
Let $u_1=(1,1)^T$, $u_2=(1,-1)^T$, $u_3=(2,0)^T$ in $\R^2$.
\textbf{(a)} Is $\{u_1,u_2,u_3\}$ a frame for $\R^2$? \textbf{(b)} If
so, find the exact frame bounds $A,B$. \textbf{(c)} Is it tight?

\smallskip
\textit{Solution.} (a) Since $u_1,u_2$ alone are linearly independent
(not parallel), they already span $\R^2$; adding $u_3$ keeps the
spanning property. By Theorem~\ref{rev:thm-spanning}, $\{u_1,u_2,u_3\}$
is a frame for $\R^2$.

(b) The synthesis matrix is $F=\begin{pmatrix}1&1&2\\1&-1&0\end{pmatrix}$.
The frame operator is
\[
  S = FF^T = \begin{pmatrix}1&1&2\\1&-1&0\end{pmatrix}
  \begin{pmatrix}1&1\\1&-1\\2&0\end{pmatrix}
  = \begin{pmatrix}1+1+4 & 1-1+0\\1-1+0&1+1+0\end{pmatrix}
  = \begin{pmatrix}6&0\\0&2\end{pmatrix}.
\]
This is already diagonal, so the eigenvalues are read off directly:
$\lambda_1=6$, $\lambda_2=2$. By Theorem~\ref{rev:thm-eigenvalues},
$B=6$ and $A=2$.

(c) Since $A\ne B$, the frame is \emph{not} tight.

\textit{Check:} the trace formula (Proposition~\ref{rev:prop-trace})
gives $\tr(S)=\norm{u_1}^2+\norm{u_2}^2+\norm{u_3}^2=2+2+4=8$, and
indeed $6+2=8$, confirming the computation.
\end{examplebox}

\begin{examplebox}[title={Review Problem 2: Canonical dual and reconstruction}]
\label{rev:problem2}
Using $\{u_1,u_2,u_3\}$ from Review Problem 1, find the canonical dual
frame, and reconstruct $x=(4,2)^T$ from its frame coefficients.

\smallskip
\textit{Solution.} Since $S=\begin{pmatrix}6&0\\0&2\end{pmatrix}$ is
diagonal, $S^{-1}=\begin{pmatrix}1/6&0\\0&1/2\end{pmatrix}$ directly.
The canonical dual vectors are
\[
  \tilde u_1 = S^{-1}u_1 = (\tfrac16,\tfrac12)^T,\quad
  \tilde u_2 = S^{-1}u_2 = (\tfrac16,-\tfrac12)^T,\quad
  \tilde u_3 = S^{-1}u_3 = (\tfrac13,0)^T.
\]
The frame coefficients of $x=(4,2)^T$ are
\[
  c_1=\inner{x}{u_1}=4+2=6,\quad c_2=\inner{x}{u_2}=4-2=2,\quad
  c_3=\inner{x}{u_3}=8.
\]
Reconstruction:
\[
  \hat x = \sum_i c_i\tilde u_i
  = 6(\tfrac16,\tfrac12)^T + 2(\tfrac16,-\tfrac12)^T + 8(\tfrac13,0)^T
  = (1,3)^T+(\tfrac13,-1)^T+(\tfrac83,0)^T = (4,2)^T = x.
\]
The reconstruction is exact, as guaranteed by
Proposition~\ref{rev:prop-canonical}.
\end{examplebox}

\begin{examplebox}[title={Review Problem 3: A tightness proof}]
\label{rev:problem3}
Let $v_1=(1,0)^T$, $v_2=(0,1)^T$, $v_3=(-1,0)^T$, $v_4=(0,-1)^T$ in
$\R^2$. Prove $\{v_i\}_{i=1}^4$ is a tight frame and find its bound.

\smallskip
\textit{Solution.} This is the $m=4$ regular polygon case of
Theorem~\ref{rev:thm-polygon} (vertices at angles $0,\pi/2,\pi,3\pi/2$),
so it is automatically tight with $A=m/2=2$. As a direct check via the
resolution-of-the-identity characterization of tightness (Chapter 4,
Proposition 4.2.1): $S=\sum_iv_iv_i^T
= 2e_1e_1^T+2e_2e_2^T=2I$, confirming $A=B=2$ without needing the
general theorem at all --- a useful sanity check technique when the
configuration is simple enough to verify directly by hand.
\end{examplebox}

\begin{examplebox}[title={Review Problem 4: The dual family, dimension and minimum}]
\label{rev:problem4}
For the frame $\{u_1,u_2,u_3\}$ of Review Problem 1, (a) find a vector
spanning $\ker F$, (b) state the dimension of the full affine family of
dual frames, and (c) confirm that the canonical dual found in Review
Problem 2 minimizes $\norm{G}_F^2$ over that family.

\smallskip
\textit{Solution.} (a) We need $c=(c_1,c_2,c_3)$ with
$c_1u_1+c_2u_2+c_3u_3=0$, i.e.\ $(c_1+c_2+2c_3,\,c_1-c_2)=(0,0)$. From
the second equation $c_1=c_2$; substituting into the first,
$2c_1+2c_3=0$, so $c_3=-c_1$. Taking $c_1=1$: $\ker F=\spn\{(1,1,-1)^T\}$.

(b) By the dimension-counting fact above (Section~\ref{sec:review-q4}),
with $n=2$ and $\dim\ker F=1$, the affine family of duals has dimension
$n\cdot1=2$ --- a plane, exactly as in Chapter~6.

(c) By Theorem~\ref{rev:thm-paraboloid} (which did not depend on the
specific frame, only on $m=n+1$), the energy function over this plane
is $\norm{\tilde F}_F^2+\norm{w}^2$, minimized uniquely at $w=0$, which
by construction is the canonical dual of Review Problem 2. No further
computation is needed: minimality is guaranteed by the general theorem.
\end{examplebox}

\section{Practice Problem Set}
\label{sec:practice-problems}

The following problems are unworked. They mirror the format of the
midterm: a mix of short computational problems and proof-style
questions. Problems marked \textbf{(C)} are intended as the take-home
computational component; all others are intended to be solvable by hand
within exam time constraints.

\begin{exercisebox}[title={Problem 1}]
Let $w_1=(3,4)^T$, $w_2=(4,-3)^T$ in $\R^2$ (note: $\norm{w_1}=\norm{w_2}=5$
and $\inner{w_1}{w_2}=0$). Without computing $S$ explicitly, determine
whether $\{w_1,w_2\}$ is tight, and if so, find $A$. (Hint: what kind of
configuration is this?)
\end{exercisebox}

\begin{exercisebox}[title={Problem 2}]
Let $\{f_i\}_{i=1}^m$ be \emph{any} frame for $\R^n$ with frame bounds
$A, B$. Prove that $\{cf_i\}_{i=1}^m$ (every vector scaled by a fixed
nonzero constant $c\in\R$) is also a frame for $\R^n$, and find its
frame bounds in terms of $A$, $B$, and $c$.
\end{exercisebox}

\begin{exercisebox}[title={Problem 3}]
Let $\{f_i\}_{i=1}^m$ be a tight frame for $\R^n$ with bound $A$, and
suppose additionally that every $f_i$ has the same norm $\norm{f_i}=r$
for all $i$. Express $A$ in terms of $m$, $n$, and $r$ only. (This
generalizes the unit-norm case, Proposition~\ref{rev:prop-trace}, to
arbitrary common norm $r$.)
\end{exercisebox}

\begin{exercisebox}[title={Problem 4}]
True or false, with justification (a one-line counterexample suffices
for any false statement):
\begin{enumerate}[label=(\alph*)]
  \item Every frame with $m=n$ vectors in $\F^n$ is an orthonormal
        basis.
  \item If $\{f_i\}_{i=1}^m$ is a tight frame for $\R^n$, then so is
        every subset $\{f_i\}_{i\in J}$ for $J\subsetneq\{1,\ldots,m\}$.
  \item The canonical dual of a Parseval frame is the frame itself.
  \item A frame operator $S$ always has $n$ distinct eigenvalues.
  \item If $A=B$ for a frame, the frame vectors must all have the same
        norm.
  \item The dimension of the affine family of dual frames for $m$
        vectors in $\R^n$ ($m>n$) equals $m-n$.
\end{enumerate}
\end{exercisebox}

\begin{exercisebox}[title={Problem 5}]
Let $\{f_i\}_{i=1}^m$ be a frame for $\R^n$ with frame operator $S$,
and let $T:\R^n\to\R^n$ be an invertible linear map. Show that
$\{Tf_i\}_{i=1}^m$ is also a frame for $\R^n$, and express its frame
operator in terms of $T$ and $S$. Under what condition on $T$ is the
new frame tight whenever the original one is?
\end{exercisebox}

\begin{exercisebox}[title={Problem 6}]
Let $\{f_1,f_2,f_3\}$ be the triangular frame of Chapters~1/3/6
($A=B=\tfrac32$). Using the parametrization $G(w)=\tilde F+wk^T$ from
Chapter~6 with $k=(1,1,1)^T/\sqrt3$, find the value of $w\in\R^2$ for
which $\norm{G(w)}_F^2=\tfrac{13}{3}$. (There are two such values by
symmetry; finding either is sufficient.) Verify your answer using
Theorem~\ref{rev:thm-paraboloid}.
\end{exercisebox}

\begin{exercisebox}[title={Problem 7 (C)}]
Computational. Generate $20$ random unit vectors in $\R^4$
(\texttt{rng.standard\_normal((4,20))}, columns normalized). Compute the
exact frame bounds $A, B$ via the frame operator's eigenvalues. Then use
gradient information (or simple coordinate-wise perturbation) to adjust
the vectors slightly to reduce the ratio $B/A$ --- you do not need a
sophisticated optimizer; even a simple loop that nudges each vector
toward better balance and re-normalizes will do. Report your starting
and ending $B/A$ ratio. (This is a hands-on preview of the frame
potential minimization in Chapter 9.)
\end{exercisebox}

\begin{exercisebox}[title={Problem 8 (C)}]
Computational. For $m=3,4,5,6,7,8,9,10$, construct the regular $m$-gon
frame in $\R^2$, compute its canonical dual, and verify numerically
that the canonical dual is simply $\tfrac1A$ times the original frame in
every case (as predicted for tight frames). Then, for each $m$, compute
$\norm{\tilde F}_F^2$ (the canonical dual's total energy) and verify it
equals $m/A$ for these unit-norm configurations (derive the predicted
formula yourself before checking it numerically).
\end{exercisebox}

\begin{exercisebox}[title={Problem 9}]
Prove: if $\{f_i\}_{i=1}^m$ is a frame for $\R^n$ with canonical dual
$\{\tilde f_i\}$, then the canonical dual of $\{\tilde f_i\}$ is
$\{f_i\}$ again (i.e.\ the canonical-dual operation is an involution).
(Hint: use Proposition~\ref{rev:prop-canonical} twice, recalling that
the frame operator of $\{\tilde f_i\}$ is $S^{-1}$.)
\end{exercisebox}

\begin{exercisebox}[title={Problem 10}]
A frame $\{f_i\}_{i=1}^m$ for $\R^n$ is called \emph{equal-norm} if
$\norm{f_i}=\norm{f_j}$ for all $i,j$. Is every tight frame equal-norm?
Is every equal-norm frame tight? Justify each answer with either a short
proof or an explicit counterexample drawn from frames appearing
somewhere in Chapters 1--6.
\end{exercisebox}

\begin{exercisebox}[title={Problem 11 (C)}]
Computational. For the frame $\{h_1,h_2,h_3\}$ ($A=1$, $B=3$), use
\texttt{scipy.optimize} to find the point $w$ on the dual-family plane
(parametrized as in Chapter~6, with the appropriate kernel vector for
this frame) at which the reconstruction-error ellipse is exactly
circular, and report both $w$ and the resulting total energy
$\norm{G(w)}_F^2$. Compare this energy to the canonical dual's minimal
energy $\tfrac43$ found in Example~5.3. (This reproduces Chapter~6, Lab
Exercise 6.2(d)--(e); if you completed that exercise already, this
problem is a direct check of your earlier work.)
\end{exercisebox}

\section{Self-Assessment Checklist}
\label{sec:checklist}

Before the midterm, you should be able to do each of the following
without consulting these notes:

\begin{itemize}[label=$\square$]
  \item State the frame definition and explain, in one sentence, why
        spanning alone suffices in finite dimensions.
  \item Given an explicit small set of vectors, determine by hand
        whether it spans, and if so compute the frame operator and its
        eigenvalues.
  \item State the relationship $S=FF^*=\sum_if_if_i^*$ and use both
        forms appropriately.
  \item Explain what it means for a frame to be tight, in three
        equivalent ways: (i) $A=B$, (ii) $S=AI$, (iii) every direction
        is ``illuminated'' equally.
  \item Prove that the vertices of a regular polygon form a tight
        frame, using either the trigonometric or roots-of-unity
        argument.
  \item Compute the canonical dual of a small explicit frame, by hand,
        via $S^{-1}$.
  \item State and use the characterization of all dual frames as
        $\tilde F + Z$ with $ZF^*=0$, and correctly count the
        \emph{dimension} of that family ($n\cdot(m-n)$, not $m-n$).
  \item Explain, in your own words, why the canonical dual is optimal
        in the minimal-norm / minimal-noise sense, and why this is
        \emph{not} the same as minimizing the anisotropy of the error
        ellipse.
  \item Sketch, for the smallest redundant case $m=n+1$, the energy
        paraboloid over the dual-family plane and the resulting error
        ellipse at an arbitrary point $w$.
  \item Translate any of the above into working NumPy code without
        referring to previous listings.
\end{itemize}